\chapter{讨论}\label{Chap:chap5}

本章讨论Swift算法的设计考量、理论分析、适用边界以及与相关工作的比较。

\section{设计考量}

\subsection{11维观测空间的设计原理}

Swift的11维观测空间设计基于信息完整性、状态可观测性、时间复杂度和扩展性四方面考量。传统强化学习拥塞控制方案（如Aurora\upcite{aurora}）通常只使用4--6个状态参数，丢失了重要的协议栈内部信息。Swift的11维观测空间包含了TCP协议栈提供的所有关键信息，包括ECN状态、拥塞控制状态机状态和拥塞事件类型。所有11个有效观测参数以及另外4项元数据通道（合计15个，嵌入OpenGym \texttt{Box}传输容器）均可从ns-3 TCP模块的标准回调接口直接获取，无需修改协议栈或引入额外的探测流量。由于观测维度固定且决策为规则式，单次状态处理只涉及常数次字段读取与比较，时间复杂度为$O(1)$；本文不报告具体的决策耗时数值，理由见第~\ref{sec:overhead}节。

\subsection{差异化窗口保留因子的选择}

Swift的差异化窗口保留因子（丢包0.70、ECN 0.75、超时0.50）的选择基于以下分析：

\textbf{丢包保留因子0.70}：传统AIMD算法在丢包时将窗口减半（保留因子0.5）。Swift选择0.70是因为在有ECN支持的网络中，适度的窗口缩减即可缓解拥塞，同时能够更快地恢复吞吐量。相比传统的0.5，保留因子0.70使得丢包后的恢复时间缩短约40\%。

\textbf{ECN保留因子0.75}：ECN标记是早期拥塞信号，此时网络可能只是轻度拥塞。保留因子0.75使窗口温和缩减，配合降窗后冻结机制释放队列压力，同时避免超时式深度回退。需要说明的是，该取值属于\emph{设计论证}而非实验结论：第~\ref{Chap:chap4}章的归档数据在丢尾队列下生成，ECN分支未被触发，因此0.75的实际效果有待在启用RED标记后的重跑中验证。

\textbf{超时保留因子0.50}：超时丢包是最严重的拥塞信号，通常意味着连续多个数据包丢失或网络路径严重中断。保留因子0.50（即窗口减半）与传统TCP的超时响应一致，确保在极端拥塞情况下能够快速释放网络资源。这是三种拥塞类型中最严厉的响应。

此外，连续缩减保护机制（最多3次连续缩减后保持窗口）的设计基于以下分析：在连续拥塞事件中，窗口经过3次保留因子为0.75的缩减后，约保留原窗口的$0.75^3 \approx 42.2\%$，已经完成显著退让。继续缩减将导致窗口过小，恢复时间过长，反而不利于整体吞吐量。

\subsection{自适应参数调优的设计}

Swift的自适应参数调优机制遵循多因素综合（最小RTT感知的RTT膨胀阈值、奖励信号EMA趋势、连续增长趋势三个维度提高鲁棒性）、渐进调整（避免剧烈变化）和边界约束（$\alpha \in [0.85, 1.30]$，基准值1.10）三个原则。

\section{理论分析}

\subsection{收敛性分析}

Swift算法的收敛性从窗口动态稳定性、拥塞响应收敛性和自适应参数稳定性三个角度进行分析。在无拥塞情况下，当窗口接近BDP时增长速度减缓；发生拥塞时，保留因子都大于0.5，窗口不会过度缩减；$\alpha$的调整受到边界约束且采用小步长，长期收敛到与网络环境匹配的值。

\subsection{公平性分析}

Swift的公平性取决于每流独立状态管理（避免流间干扰）、差异化响应一致性（所有Swift流对相同拥塞信号采用相同保留因子）和与传统算法共存的兼容性。

\subsection{复杂度分析}

\textbf{时间复杂度}：单次决策的时间复杂度为$O(1)$，包括观测空间处理、拥塞检测、参数调优和窗口计算。

\textbf{空间复杂度}：每流状态由常量规模字段构成，包括两个长度40的滑动窗口（投递速率过滤窗口与RTT采样窗口）、一个容量上限4096的投递速率采样队列（其实际占用随路径RTT与确认密度自适应收缩）、BDP估计值、$\alpha$以及若干计数器，因此单流为$O(1)$，总空间复杂度为$O(n)$（$n$为并发流数量）。

\textbf{通信开销}：每次回调产生一次观测-动作交互，消息大小约120字节（15个64位整数）。

\section{适用边界与改进方向}
\label{sec:limitation-future}

\subsection{缓冲区相对BDP过深场景中的延迟抬升}

第~\ref{Chap:chap4}章的结果显示，Swift在十个场景中的平均时延高于NewReno/CUBIC：S8（\texttt{dc\_100m}）$+123.3$\%、S13（\texttt{wifi\_n}）$+100.4$\%、S18（\texttt{lte\_poor}）$+71.6$\%、S14（\texttt{wifi\_legacy}）$+61.6$\%、S16（\texttt{nr\_5g\_edge}）$+60.0$\%、S7（\texttt{dc\_500m}）$+57.0$\%、S17（\texttt{lte\_good}）$+33.7$\%、S4/S5约$+22.7$\%、S6（\texttt{symmetric\_low}）$+21.3$\%。

这些场景具有共同的结构特征，据此可以给出比"$\alpha$上界偏大"更准确的机理判断。仿真器按$Q = \max(BDP \times n / MTU, 100)$个数据包为瓶颈队列定尺，在低带宽、短传播时延的路径上BDP极小，$Q$因而被100个数据包的下界保护决定。以S8为例：瓶颈100\,Mbps、基线RTT 18\,$\mu$s，BDP仅约225字节，而100个数据包的缓冲区容量达150\,kB，二者相差近三个数量级。与此同时，Swift的最终窗口被夹紧在$[4 \times MSS, \max(4 \times BDP, 200 \times MSS)]$，当$4 \times BDP$远小于$200 \times MSS = 288$\,kB时，上界由后者决定。于是出现如下情形：\emph{目标窗口机制$\alpha \times BDP$虽然计算正确，但其结果远低于窗口下界与安全上界所构成的可行区间，窗口实际由兜底常量而非BDP目标主导}。此时无论$\alpha$取$0.85$还是$1.30$，窗口都稳态驻留在队列中。

这一诊断把改进方向从"调小$\alpha$"修正为"消除兜底常量与BDP估计之间的量级失配"，可进一步探索的方案包括：（1）将$200 \times MSS$的兜底上界改为仅在BDP估计\emph{尚未建立}时生效，一旦估计有效即完全交由$\max(k \times BDP, W_{min})$约束；（2）引入排队延迟显式估计$q_{est} = RTT_{smoothed} - RTT_{min}$作为窗口上界的直接约束，使窗口在排队持续增长时无条件回落，而不依赖BDP估计的绝对精度；（3）采用多时间尺度BDP估计，分离传播时延与排队时延的贡献，以提高低BDP路径上估计的相对精度。需要说明的是，上述三项均属设计推论，其效果需在完成全矩阵重跑后由实验判定。

\subsection{跨域广域网场景的吞吐量回落}

在S9（\texttt{cross\_dc\_wan}，1\,Gbps瓶颈、基线RTT 10.04\,ms）中，Swift聚合吞吐量709.7\,Mbps，相对NewReno/CUBIC均值回落28.97\%。同一场景中Swift的时延为四者最低（5.04\,ms）且丢包为零，说明这是窗口偏小而非拥塞失控。

该现象的机理已被定位，且不属于参数取值问题。生成本文数据的实现修订采用$DR = segmentsAcked \times segSize / lastRtt$计算投递速率样本：分子仅为单次确认的字节数，分母却是整个RTT，因此估计值被系统性压低约等于每RTT确认个数的倍数。在1\,Gbps瓶颈、10\,ms基线RTT的路径上该倍数可达数百，BDP估计随之崩塌，$\alpha \times BDP$目标窗口长期低于实际需求，窗口被压在$200 \times MSS$的兜底上界附近，吞吐量被锁定在约$200 \times MSS / RTT$的水平。第~\ref{Chap:chap3}章第3.6.1节给出了该推导的完整形式。

v3.0.0已将该估计器替换为时间窗口差分形式$DR = [B(t) - B(t-\Delta)]/\Delta$，其中$\Delta \geq 2 \times RTT_{min}$，分子与分母取自同一时间区间，不再依赖每RTT确认个数。因此本项限制预期在重跑后消失或显著缓解，但\emph{在获得重跑数据之前，本文不宣称该问题已被解决}。除此之外，长RTT路径上仍值得探索的方向包括：细化长RTT条件下$\alpha$的更新步长，以及在有安全约束的前提下引入更积极的带宽探测。

\subsection{实验证据的范围限制}

本文的定量结论受两项范围限制约束，二者都不影响机理分析的有效性，但限定了结论的适用范围。

第一，\textbf{单一随机种子}。\path{logs/}下的产物每个(场景, 协议, 设置)三元组仅对应一次仿真运行，因此本文不报告跨种子的标准差与置信区间，并据此仅把量级明确、机理可解释的差异写入结论（详见第~\ref{Chap:chap4}章第4.1.5节）。ns-3在给定\texttt{RngRun}下完全确定，故现有结果可精确复现；其对随机实现的稳健性需借助\texttt{main.py --num-seeds}完成多种子重跑后确认。

第二，\textbf{实现修订边界}。全部数据生成于v3.0.0修订之前，该修订变更了投递速率估计器、启用了RED的ECN标记通路与全协议对称ECN、修正了ECE/CWR事件的信号归类，并把奖励自适应改为基线相对判据。因此现有数据无法单独分离出这四项变更各自的贡献。完成全矩阵重跑既是消除第一项限制的手段，也是量化第二项变更效果的前提，因而是后续工作中优先级最高的一项。

\subsection{参数自动调优}

当前Swift的保留因子（0.70/0.75/0.50）、$\alpha$边界（$[0.85,1.30]$）、连续递减阈值（$D_{max}=3$）和降窗后冻结窗口增长的ACK事件数（4）为当前实现中的手动调优预设值。后续研究可通过元学习\upcite{meta_learning}或贝叶斯优化探索参数的在线自动搜索，但相关结论需要在完整实验矩阵重跑后再报告。

\subsection{深度强化学习增强}

Swift当前采用基于规则的策略，其可解释性强但表示能力有限。后续研究可探索离线强化学习训练神经网络策略，再通过模型蒸馏提取为紧凑规则，兼顾表示能力和实时性；同时需要引入安全强化学习约束，确保策略不会产生危险的窗口调整。

\subsection{真实部署验证}

需将算法实现为Linux内核TCP模块或QUIC用户态协议栈，在覆盖近端高速、蜂窝/无线接入、广域网与卫星链路的真实或半实物测试床上验证。仿真与真实环境的差异（如NIC offload、中断合并、NUMA效应、硬件时间戳精度、终端操作系统调度和无线链路波动）可能影响实际性能。QUIC协议的用户态实现可简化部署流程，且QUIC的多流特性可实现更精细的流间调度。

\section{与相关工作的比较}

\subsection{与Aurora的比较}

Aurora\upcite{aurora}是基于深度强化学习的拥塞控制算法，与Swift的主要区别包括：

\begin{table}[htbp]
\centering
\caption{Swift与Aurora的详细对比}
\label{tab:vs_aurora}
\begin{tabular}{|l|c|c|}
\hline
\textbf{特性} & \textbf{Swift} & \textbf{Aurora} \\
\hline
有效观测维度 & 11（另含4项元数据） & 10 \\
ECN状态利用 & 是 & 否 \\
CA状态利用 & 是 & 否 \\
决策机制 & 基于规则 & 深度神经网络 \\
可解释性 & 高 & 低 \\
训练需求 & 无 & 需要大量训练 \\
决策复杂度 & $O(1)$ & 神经网络前向推理 \\
参数调优 & 在线自适应 & 固定（训练后） \\
\hline
\end{tabular}
\end{table}

\subsection{与DCTCP的比较}

DCTCP\upcite{dctcp}是利用ECN进行细粒度调整的拥塞控制算法。

\begin{table}[htbp]
\centering
\caption{Swift与DCTCP的详细对比}
\label{tab:vs_dctcp}
\begin{tabular}{|l|c|c|}
\hline
\textbf{特性} & \textbf{Swift} & \textbf{DCTCP} \\
\hline
ECN利用方式 & 状态直接检测 & 标记比例连续调整 \\
拥塞响应 & 差异化保留因子 & 统一公式 \\
丢包响应保留因子 & 0.70 & 0.50 \\
ECN响应保留因子 & 0.75 & 可变（基于$\alpha$） \\
参数自适应 & 是 & 否 \\
协议栈状态利用 & 是 & 否 \\
多信号融合 & 是 & 否 \\
\hline
\end{tabular}
\end{table}

\subsection{与BBR的比较}

BBR\upcite{bbr}是基于带宽和延迟估计的拥塞控制算法。

\begin{table}[htbp]
\centering
\caption{Swift与BBR的详细对比}
\label{tab:vs_bbr}
\begin{tabular}{|l|c|c|}
\hline
\textbf{特性} & \textbf{Swift} & \textbf{BBR} \\
\hline
主要拥塞信号 & 丢包+ECN+CA状态 & RTT测量 \\
带宽探测方式 & 自适应参数调优 & 周期性增益调整 \\
丢包响应 & 是 & 否（设计上） \\
ECN利用 & 是 & 否 \\
控制模式 & 窗口控制 & 速率控制 \\
低RTT场景带宽利用率 & 接近饱和 & 场景相关 \\
中速场景带宽利用率 & 较高 & 较高 \\
与CUBIC共存公平性 & 场景相关，依赖队列与RTT条件 & 场景相关，依赖增益周期 \\
\hline
\end{tabular}
\end{table}

\subsection{与PCC的比较}

PCC\upcite{pcc}采用在线实验探索最优速率，而Swift采用预设规则结合自适应参数调优。PCC需要在线进行速率探索实验，可能导致短期性能下降；Swift不需要显式探索，性能更稳定。

\section{敏感性分析}

参数敏感性的理论分析应服务于对实现机制的解释。当前实现采用固定边界与事件驱动更新逻辑：ECN保留因子为0.75，普通丢包保留因子为0.70，超时保留因子为0.50；乘性增加因子$\alpha$在$[0.85,1.30]$范围内根据RTT膨胀、奖励信号EMA和连续增长趋势动态调整。本节不引入未由日志支撑的额外参数扫描结果，而从机制角度说明这些取值如何影响稳定性与吞吐量。

\subsection{窗口保留因子的理论依据}

丢包保留因子0.70的选择基于收敛速度与吞吐量的权衡：在AIMD框架下，保留因子$\beta$决定了窗口从$W_{max}$恢复到$W_{max}$所需的时间$T_{recover} \propto \frac{1-\beta}{\gamma}$个RTT。$\beta=0.70$相比传统$\beta=0.50$使恢复时间缩短40\%，但丢包后的即时退让幅度也相应减小。在有ECN支持的网络中，丢包事件频率显著降低，因此更温和的丢包响应是合理的。

ECN保留因子0.75的选择基于ECN信号的语义：ECN标记表示队列深度超过阈值但尚未溢出，属于早期拥塞信号。相较普通丢包响应，0.75保留因子能够在发生早期拥塞通知时更温和地释放队列压力；相较超时响应，它又避免了不必要的窗口深度回退。该参数与降窗后冻结机制共同作用，用于降低连续ACK事件中立即反弹造成的二次排队风险。

\subsection{连续递减保护阈值的数学分析}

在连续ECN标记事件中，窗口经过$D_{max}$次保留因子为$\beta_{ecn}=0.75$的收缩后，剩余窗口比例为$0.75^{D_{max}}$。当$D_{max}=3$时，算法在连续拥塞信号下已经完成显著退让；继续无约束缩减容易造成窗口过低和恢复迟滞。因此实现中采用连续递减保护：当连续缩减计数超过阈值后暂时保持当前窗口，同时通过降窗后冻结若干ACK事件抑制立即增窗，使队列有时间排空。

\subsection{RTT膨胀比阈值的设计考量}

RTT膨胀调节采用路径感知的动态阈值，而非固定阈值。实现首先根据最小RTT构造松弛量$s = 0.3 + 0.15 \sqrt{\max(0, RTT_{min,ms}-1)}$，再形成$1+s$、$1+2s$和$1+4s$三个阈值。短RTT路径对排队更敏感，阈值相对更紧；长RTT路径允许更大的相对抖动，阈值相对放宽。该设计避免在远距离路径上因正常RTT波动过早收紧窗口，也避免在低RTT路径上放任队列持续积累。

\section{计算开销分析}
\label{sec:overhead}

Swift当前实现采用事件驱动的规则决策与每流独立状态表。由于本仓库日志未包含跨硬件平台的微基准测量，本节不报告未验证的决策延迟、内存占用或CPU利用率数值，而从代码结构给出可复核的复杂度分析。

\subsection{时间复杂度}

每次ACK或拥塞事件触发时，算法只执行常数次状态读取、拥塞类型判定、$\alpha$更新和窗口计算。带宽估计使用固定长度为40的滑动窗口，RTT采样窗口长度同样为40，且最大交付速率由窗口队列维护。因此，单次决策的时间复杂度为$O(1)$，不会随仿真总时长或历史样本数量线性增长。

\subsection{空间复杂度}

每个TCP连接维护独立状态，包括长度为40的交付速率窗口、长度为40的RTT采样窗口、BDP估计值、$\alpha$、连续增长/缩减计数器、丢包/ECN计数器和降窗后冻结计数器等常量规模字段。因此，单流空间复杂度为$O(1)$，总空间复杂度随并发连接数线性增长，即$O(N)$。

\subsection{工程开销来源}

主要开销来自OpenGym/ZMQ同步交互、Python智能体决策和ns-3回调状态采集。当前实现优先保证强化学习动作与离散事件仿真时序一致，因此未在OpenGym/ZMQ路径上启用MTP多线程并行仿真。若未来进行内核态或用户态协议栈移植，应重新测量实际平台上的决策时延和资源占用。

\section{理论分析深入探讨}

\subsection{多信号融合的信息论基础}

多信号融合拥塞检测的设计可以从信息论角度进行分析。设网络的真实拥塞状态为随机变量$C \in \{0, 1\}$，观测到的信号集合为$\mathbf{S} = \{S_1, S_2, ..., S_k\}$。根据贝叶斯定理：
\begin{equation}
P(C=1|\mathbf{S}) = \frac{P(\mathbf{S}|C=1) \cdot P(C=1)}{P(\mathbf{S})}
\end{equation}

Swift的分层优先级机制可以理解为对不同拥塞信号可靠性的排序：慢启动阈值回调对应的丢包/超时信号具有最高优先级；ECN CE或ECE状态表示队列已经触发显式拥塞通知，具有次级优先级；RTT膨胀和恢复状态只参与参数调节，不作为独立降窗触发条件。这样的排序与实现中的\texttt{\_detect\_congestion}函数保持一致。

\subsection{差异化响应的最优控制分析}

差异化窗口保留因子的设计可以从最优控制理论角度分析。优化目标是最大化长期吞吐量同时最小化延迟：
\begin{equation}
J = \sum_{t=0}^{T} \left[ -\alpha \cdot throughput(t) + \beta \cdot delay(t) + \gamma \cdot |u(t)|^2 \right]
\end{equation}

Swift的差异化保留因子（丢包0.70、ECN 0.75、超时0.50）可以理解为对不同拥塞严重程度的分层近似。

\subsection{自适应参数调优的强化学习视角}

自适应参数调优机制可以形式化为强化学习问题。Swift当前的基于规则策略可以表示为：
\begin{equation}
\pi(s_t) = \begin{cases}
+\Delta_{\text{up}} & \text{如果~} r < 1+s \text{，表示排队膨胀较低} \\
+\Delta_{\text{small}} & \text{如果~} 1+s \le r < 1+2s \text{，表示温和增长} \\
-\Delta_{\text{down}} & \text{如果~} r > 1+4s \text{，表示重度排队} \\
+\Delta_{\text{ema}} & \text{如果~} R_{fast} > R_{base} + m \\
-\Delta_{\text{ema}} & \text{如果~} R_{fast} < R_{base} - 4m \\
+\Delta_{\text{stable}} & \text{如果~} consecutive\_growth > 8 \\
0 & \text{否则}
\end{cases}
\end{equation}
其中$s = 0.3 + 0.15\sqrt{\max(0, RTT_{min,ms}-1)}$，$r=RTT_{last}/RTT_{min}$，$R_{fast}$与$R_{base}$分别为奖励信号的快线（$\eta_f=0.15$）与慢线基线（$\eta_b=0.02$）指数移动平均，$m = \max(0.25, 0.1|R_{base}|)$为自适应容差。该形式体现了当前实现中的两项关键设计：路径感知的动态RTT阈值，以及以基线相对偏离而非绝对水平判断奖励趋势——后者的必要性在于环境侧奖励在绝大多数ACK事件上为正，绝对阈值会使调节退化为单向棘轮。

\subsection{收敛性的Lyapunov分析}

定义Lyapunov函数：
\begin{equation}
V(t) = \frac{1}{2}(cwnd(t) - cwnd^*)^2 + \frac{1}{2}(queue(t) - queue^*)^2
\end{equation}

当$cwnd > cwnd^*$且出现ECN拥塞通知时，当前实现按$cwnd(t+1) = 0.75 \cdot cwnd(t)$进行窗口退让；若出现普通丢包或超时，则分别按0.70或0.50执行更强响应。该分层缩减机制有助于使窗口向目标区间回落。

\subsection{公平性的博弈论分析}

在多流场景下，各流的拥塞控制可以建模为非合作博弈。每个流$i$的效用函数：
\begin{equation}
U_i(x_i, x_{-i}) = \log(x_i) - \beta \cdot delay(X)
\end{equation}

第~\ref{Chap:chap4}章的测量结果支持该分析：Swift在19个纯TCP场景中的18个取得Jain指数1.000、余下一个为0.999，均值为四种协议中最高（NewReno 0.988、CUBIC 0.993、BBR 0.812）。其结构性原因在于Swift的目标窗口由路径的BDP估计决定而非由丢包事件决定：共享同一瓶颈的三条流最终会收敛到相近的BDP估计，因而目标窗口天然对称；而基于丢包的算法依赖丢包在流间的随机分布来分配带宽，该分布的方差随RTT增长而放大，这解释了NewReno与CUBIC在GEO卫星场景（S19）中公平性分别降至0.823与0.912的现象。

\section{本章小结}

本章讨论了Swift算法的设计考量、理论分析、适用边界以及与相关工作的比较。11维观测空间基于信息完整性和可观测性原则设计；差异化保留因子基于拥塞信号语义选择；路径感知$\alpha$调节、时间窗口投递速率与BDP窗口化最大值滤波、降窗后冻结和连续缩减保护共同构成当前实现的稳定性机制；测量得到的Jain公平性均值1.000为四种协议最高，其结构性原因是目标窗口由BDP估计而非丢包事件的随机分布决定。

关于适用边界，本章给出了两项比"参数需要调优"更具体的诊断。第一，十个低速率场景中的时延抬升（$+21.3$\%$\sim+123.3$\%）源于结构性量级失配：当路径BDP远小于$200 \times MSS$的兜底上界时，窗口实际由该常量而非$\alpha \times BDP$目标主导，因此改进方向应是使兜底约束与BDP估计相协调，而非调小$\alpha$。第二，跨域WAN场景的吞吐回落（$-28.97$\%）源于生成数据时所用的单次确认投递速率公式对带宽的系统性低估，v3.0.0已将其替换为时间窗口差分形式，但在完成重跑之前不宣称该问题已解决。此外，本章明确了两项证据范围限制——单一随机种子与实现修订边界——并据此把全矩阵多种子重跑列为后续工作的最高优先级。与Aurora、DCTCP、BBR等相关工作相比，Swift在多信号融合、差异化拥塞响应和自适应参数调优方面具有差异化贡献；计算开销分析表明单次规则决策与每流状态维护均为常量规模，主要工程开销来自OpenGym/ZMQ同步交互。
